% =============================================================================
%  Section 2 — Background and Motivation
% =============================================================================
\section{Background and Motivation}
\label{sec:bg}

This section establishes the conceptual baseline of multi-tenant LLM agent
runtimes (Section~\ref{sec:bg-kernel}), exposes the limitations of conventional,
layer-separated isolation (Section~\ref{sec:bg-lim}), characterizes the novel attack
surfaces introduced by LLM agents (Section~\ref{sec:bg-attack}), and derives the
design goals that motivate Neuron (Section~\ref{sec:bg-goals}).

% -----------------------------------------------------------------------------
\subsection{Multi-Tenant LLM Agent Runtimes}
\label{sec:bg-kernel}

A multi-tenant LLM agent runtime extends the classical process abstraction
to agent workloads. Where a conventional process binds a program counter and a
file-descriptor table to a scheduling entity, an agent process binds a
\emph{turn}---an atomic interaction comprising a prompt, a sequence of tool
calls, and an LLM-mediated control flow---to a runtime-managed execution
context. The runtime is responsible for scheduling turns, mediating tool
execution, accounting for resource consumption (principally LLM tokens, which
play the role of CPU cycles), and enforcing isolation among co-resident
agents.

Three properties of modern agent workloads distinguish them from conventional
processes and elevate them to first-class runtime concerns:

\begin{itemize}
    \item \textbf{Long-running and stateful.} Agents persist memory, plans, and
    intermediate artifacts across turns and sessions, creating durable state
    that must be isolated and accounted for over long horizons.
    \item \textbf{Dynamic and self-spawning.} An agent may, in the course of a
    single turn, spawn sub-agents to delegate sub-tasks. The resulting spawn
    tree is constructed at runtime from natural-language decisions, not from a
    static process graph.
    \item \textbf{Multi-tenant.} A single runtime instance concurrently serves
    multiple tenants that share low-level substrates---an event bus for
    inter-agent communication, a virtual file system (VFS) for artifact
    storage, and a pool of virtual threads for execution---while requiring
    strict isolation and per-tenant fairness.
\end{itemize}

Existing frameworks address these properties at the application layer.
AutoGen~\cite{wu2023autogen} and LangGraph~\cite{langgraph2024} provide
orchestration primitives but delegate resource accounting to user-space token
counters and access control to static policy files. The Model Context
Protocol~\cite{anthropic2024mcp} standardizes tool exposure yet is silent on
resource governance. In all cases, the governance machinery lives above the
runtime boundary, observing the workload from the outside.

% -----------------------------------------------------------------------------
\subsection{Limitations of Layer-Separated Isolation}
\label{sec:bg-lim}

Classical isolation designs separate concerns across three largely independent
layers, each with its own mechanism and its own audit trail:

\begin{itemize}
    \item \textbf{Resource control} (quota managers such as Linux
    cgroup~v2~\cite{linux2024cgroupv2}) accounts for and caps resource
    consumption---CPU, memory, and, in our setting, LLM tokens. Resource
    decisions are recorded in quota-manager-local logs.
    \item \textbf{Execution isolation} (sandboxes, seccomp, containers such as
    gVisor~\cite{young2023gvisor} and Firecracker~\cite{agache2020firecracker})
    confines what an execution context can do at the system-call level.
    Isolation decisions are recorded in sandbox-local logs.
    \item \textbf{Access control} (capabilities~\cite{mercer2012capsicum,klein2009sel4},
    declarative policies) determines whether an operation is
    permitted. Permission decisions are recorded in permission-local logs.
\end{itemize}

This separation is benign when the three layers are genuinely independent. It
becomes \emph{unsafe} under multi-tenant agent workloads because the layers
share a common resource pool and a common execution substrate. Two coupled
pathways, invisible to any single layer, undermine the entire design:

\paragraph{Security starvation under resource contention.}
The permission layer is not a free abstraction; it consumes CPU to evaluate
rules, memory to retain audit state, and locks to coordinate concurrent
decisions. When a malicious tenant mounts a resource-contention attack---for
example, flooding the event bus with broadcasts or hammering the VFS with
writes---the security module is itself a victim of the contention it is
supposed to police. Under sufficient pressure, policy evaluation is delayed,
audit buffers overflow, and permission checks degrade to permissive fallbacks.
The security layer thus fails \emph{because of} a resource condition that it
has no mechanism to observe, let alone act upon.

\paragraph{Privilege escalation in contention windows.}
The converse pathway is equally dangerous. Between the moment resource
pressure begins to build and the moment the resource layer trips a hard
quota-exceeded signal, there exists a transient \emph{contention window} in
which the runtime is impaired but not yet in failure. During this window, an
attacker can issue a privileged operation---a cross-tenant VFS access, a
destructive tool call---that the permission layer, still nominally functional,
admits because it has no visibility into the resource state and no causal link
to the ongoing attack. By the time either layer logs the event, the two
records sit in disjoint trails with no shared identifier to correlate them.

The root cause is the absence of a \emph{shared causal identifier}.
Without a single trace that threads through quota-manager, sandbox, and
permission decisions, cross-layer correlation is impossible: the runtime
cannot reconstruct that ``the permission grant at time $t$ occurred inside the
contention window opened by the broadcast flood at time $t-\delta$.'' Each
layer operates---and fails---in isolation.

% -----------------------------------------------------------------------------
\subsection{Novel Attack Surfaces in LLM Agents}
\label{sec:bg-attack}

Beyond the coupled-failure mode, LLM agents introduce attack surfaces that
classical isolation models were not designed to handle, because they violate
two foundational assumptions: that privileges bind at start-up, and that
escalation proceeds through explicit system calls.

\paragraph{Spawn-time privilege escalation.}
In a classical system, a child process inherits a subset of its parent's
privileges; a downgraded parent cannot produce an upgraded child. In an agent
runtime, a sub-agent is constructed by instantiating a fresh execution context
through a tool call. If the spawn path does not explicitly propagate the
parent's effective privilege profile, the child defaults to a permissive
baseline---\emph{spawn itself becomes an escalation}. A parent agent that has
been deliberately downgraded (for example, a reviewer restricted to read-only
tools via a \texttt{*:deny} profile) can thereby produce a child with full
default privileges, defeating the downgrade in a single tool call.

\paragraph{Tenant-identity loss across spawn boundaries.}
Multi-tenant isolation depends on every execution context carrying a tenant
identifier that the access-control layer can check. If the spawn path neglects
to propagate the tenant identity to the child, the child's tenant field is
null, and tenant-ownership checks---which conservatively skip when the
caller's tenant is unknown---silently permit cross-tenant access. A sub-agent
spawned by a tenant-$A$ parent can then read or mutate tenant-$B$ artifacts.

\paragraph{Natural-language privilege escalation without spawn-tree context.}
An agent requests permission to invoke a sensitive tool through a structured
approval channel. If the request payload carries only the agent identifier,
the tool name, and the target---but not the position of the requesting agent
within the spawn tree---then a human approver (or an automated policy) cannot
distinguish a top-level request from one originating at depth three. A deeply
nested sub-agent, precisely the context in which an attacker is most likely to
attempt social-engineering escalation, is indistinguishable from a benign
top-level call. The approval channel, designed for flat process models, becomes
the weak link.

\paragraph{Resource-contention flooding.}
Finally, the shared substrates themselves offer an unthrottled attack surface.
Without per-tenant rate limits at the resource layer, a malicious agent can
sustain high-frequency event-bus broadcasts and VFS writes, consuming SSE
fan-out capacity and lock contention budget \emph{before} the permission layer
has any opportunity to intervene. The attack succeeds entirely within the
layer-separated blind spot.

% -----------------------------------------------------------------------------
\subsection{Governance Gap and Design Goals}
\label{sec:bg-goals}

The foregoing analysis reveals a single governance gap with two facets: the
layers are \emph{causally disconnected} (no shared identifier), and the spawn
boundary is \emph{semantically unguarded} (privileges and tenant identity are
not enforced to be non-increasing). From this gap we derive the design goals
that Neuron must satisfy:

\begin{description}
    \item[G1. Unified causal audit.] Every governance decision---resource,
    isolation, or permission---must carry a single per-turn trace identifier
    and be aggregated into a unified audit chain, so that cross-layer
    correlation and real-time attack truncation become possible.

    \item[G2. Spawn-time privilege non-increase.] A child agent's privileges
    must be a subset of its parent's; spawn must never constitute an
    escalation, closing the attack surface of Section~\ref{sec:bg-attack}.

    \item[G3. Tenant isolation across spawn.] Tenant identity must propagate
    across spawn boundaries, and tenant-ownership checks must execute as a
    hard pre-check that cannot be bypassed by permissive modes.

    \item[G4. Depth-aware escalation.] Every approval request must carry the
    spawn-tree context (depth and parent chain), and destructive operations
    originating from deeply nested sub-agents must be denied automatically.

    \item[G5. Resource-layer throttling.] The shared substrates must enforce
    per-tenant rate limits that act \emph{before} the permission layer, so
    that contention is damped at the source rather than policed after damage.

    \item[G6. Zero regression.] All closures must be transparent to headless
    execution, preserving existing fallback contracts so that the runtime's
    correctness test suite remains green.
\end{description}

Goals G1 and G5 target the coupled-failure mode of Section~\ref{sec:bg-lim};
goals G2--G4 target the LLM-specific attack surfaces of
Section~\ref{sec:bg-attack}; G6 constrains the design space to evolvable,
non-disruptive closures. The Neuron architecture presented in
Section~\ref{sec:arch} is designed to satisfy these goals jointly, rather than
in the layer-separated fashion that motivates this work.
